% ============================================================
% Gini Impurity vs. Entropy Comparison Table
% CSC491 Textbook, Chapter: Machine Learning
%
% Preamble requirements:
%     \usepackage{booktabs, array, xcolor}
% ============================================================

\begin{table}[ht]
\centering
\setlength{\tabcolsep}{10pt}
\renewcommand{\arraystretch}{1.7}
\caption{%
  scenarios where one metric may be more practical
}
\label{tab:gini-vs-entropy}
\begin{tabular}{
  >{\raggedright\arraybackslash}p{3.0cm}
  >{\raggedright\arraybackslash}p{5.0cm}
  >{\raggedright\arraybackslash}p{5.0cm}
}
\toprule
\textbf{Scenario}
  & \textbf{Gini Impurity}
  & \textbf{Entropy (Information Gain)} \\
\midrule
Split Behavior
  & Favors dominant classes; tends to isolate the largest group first
  & Produces more balanced partitions across classes \\

Dataset Size
  & Efficient for large, high-dimensional datasets
  & Better suited to structured datasets with balanced classes \\

Sensitivity
  & Less sensitive to small probability differences
  & sensitive to subtle shifts in class distribution \\

\bottomrule
\end{tabular}

\smallskip
\footnotesize
\textit{Rule of thumb:} When in doubt, start with Gini, it is faster and the default
in most production libraries. Switch to entropy if your classes are well-balanced and
you want the tree to weight minority classes more carefully.

\end{table}